\section{Discussion}

\subsection{Security Tool or Harness?}

ActPlane can express security-relevant policies, such as keeping secret-derived
data local. However, its primary target is an agent harness: a deterministic
operating contract for a cooperative but forgetful agent. This affects design
choices. A classic reference monitor would prioritize pre-operation denial for
all security claims. ActPlane also needs audit-only guidance, harness-level
termination when pre-operation arguments are unavailable, and remediation text
that lets the agent complete the original task through an allowed path.

\subsection{Where the Claim Must Stay Honest}

The strongest novelty claim is not the taint algorithm. In-kernel label
propagation and provenance enforcement are established ideas. ActPlane must be
evaluated as a systems contribution that combines those ideas with a new
workload and interface: autonomous agents with rich tool access, natural-language
policy sources, and corrective feedback.

The final paper should be careful about three boundaries.

\textbf{Block is LSM-only.} Tracepoint observation cannot be described as
\texttt{block}. The current implementation correctly ignores \texttt{block}
rules in tracepoint mode and warns the operator.

\textbf{Prompt files are aspirational.} A corpus of instruction files shows what
developers want agents to do. It does not by itself prove how often agents
violate those rules. That is the job of the agent evaluation.

\textbf{Coarse taint can over-approximate.} Object-level labels are practical,
but they may carry more provenance than a byte-level system would. The evaluation
must measure false positives and demonstrate useful declassification and gate
patterns.

\section{Conclusion}

ActPlane moves agent harness enforcement below the tool layer. It compiles
agent-oriented information-flow policies into eBPF/BPF-LSM checks, tracks labels
across process, file, and endpoint objects, and returns kernel-detected
violations as corrective feedback to the agent. The evaluation plan is designed
to test bypass coverage, feedback-driven recovery, false positives, overhead,
and coverage over real agent instruction files without reporting numbers before
they can be reproduced.
